\par
Faithful transport preserves the uncertainty an estimator actually supplies.
It does not establish that the untreated trend restriction is credible. We
now turn from conversion and covariance contracts to the effect of weakening
that restriction, keeping uncertainty about the observed pretrend in the
calculation rather than treating its estimate as known.

\section{How departures from parallel trends change the conclusion}
\label{exec22:hd-scope}

A funding manager evaluating a branch offer needs to distinguish a plausible
departure from parallel trends from uncertainty about the preperiod trend itself.
HonestDiD is a useful reference for this distinction. The selected R revision
declares version 0.2.8 and implements sensitivity inference associated with
Rambachan and Roth. The paper supplies identification and inferential arguments;
the source supplies a particular numerical interface. Neither makes an invalid
event-study design valid simply because a confidence interval was returned.
The proposed library should retain the method family, while requiring a separate
qualification of the R worker, its solvers and its conversion of bank targets
before operational adoption~\citep{rambachan2023,exec22_honestdid_engine}.

\subsection{Two restrictions with different zero points}
\label{exec22:hd-restrictions}

Consider an offer that starts at time $1$, with one earlier comparison at $-1$
and an untreated reference period $0$. Let $\delta_t$ be the difference between
treated and comparison groups' untreated mean outcomes, expressed relative to
its value at time $0$, so $\delta_0=0$. Assume no anticipation through time $0$,
a fixed population and a coherent, equally spaced calendar. Write $x=\delta_{-1}$
and $y=\tau+\delta_1$, where $\tau$ is the postperiod treatment effect in the same
outcome units. These are population contrasts, not estimated coefficients.
An outcome can be quarterly contribution in dollars per eligible exposure,
provided the same denominator and accounting convention apply in all periods.

A relative-magnitude restriction with multiplier $m\geq0$ bounds the next
untreated first difference by $m$ times the absolute preperiod first difference.
An absolute smoothness restriction with $M\geq0$ bounds the change in that
first difference. In this three-period case they give
\begin{align}
 |\delta_1|&\leq m|x|,
 &\mathcal I_{\rm RM}(x,y;m)&=[y-m|x|,\ y+m|x|],
 \label{exec22:hd-rm}\\
 |\delta_1+x|&\leq M,
 &\mathcal I_{\rm SD}(x,y;M)&=[y+x-M,\ y+x+M].
 \label{exec22:hd-sd}
\end{align}
The first interval follows by substituting every admissible $\delta_1$ into
$\tau=y-\delta_1$; the second follows from
$\delta_1\in[-x-M,-x+M]$. Every endpoint is attainable under the stated
restrictions. With $x=-2$ and $y=4$, zero relative magnitude gives effect $4$,
whereas zero smoothness gives effect $2$. The latter continues a linear
untreated differential trend. It does not impose a flat trend. With $m=M=1$
the intervals are $[2,6]$ and $[1,3]$, but the numerical equality of the tuning
values has no substantive meaning: $m$ is dimensionless, while $M$ is an
absolute change in the outcome difference on the chosen discrete calendar.
Changing dollars to cents multiplies $M$ and the target by $100$ and leaves
$m$ unchanged. Changing the spacing or aggregation of periods requires a new
restriction argument; it cannot be handled by relabelling an axis
\citep{rambachan2023,exec22_honestdid_reference}.

The earlier curvature analysis develops the multiperiod smoothness kernel.
The present comparison adds a separate concern: setting an estimated pretrend
equal to its population value can remove precisely the uncertainty that a
sensitivity procedure is supposed to propagate. An estimated $\widehat x=0$
does not establish $x=0$. Failing to reject a pretrend test does not establish
that equality either. A proposed workflow must pass the preperiod estimates
and their joint covariance into its sensitivity calculation, alongside a
substantive justification of the restriction.

\subsection{A transparent projection reference with noisy preperiods}
\label{exec22:hd-projection}

For a reproducible design reference, stipulate the exact observation law
\begin{equation}
 \widehat\beta=(\widehat x,\widehat y)'\sim
 N_2\bigl((x,y)',\Sigma\bigr),\qquad \Sigma\succ0,
 \label{exec22:hd-law}
\end{equation}
with known covariance. This is an authored Gaussian experiment, not an
assertion about a fitted banking panel. For $0<\alpha<1$ set
$r^2=\chi^2_{2,1-\alpha}=-2\log\alpha$. The ellipse
\begin{equation}
 E=\{\beta:(\beta-\widehat\beta)'\Sigma^{-1}
                  (\beta-\widehat\beta)\leq r^2\}
 \label{exec22:hd-ellipse}
\end{equation}
contains the true vector with probability $1-\alpha$. Project all the
relative-magnitude target intervals over this ellipse. For $s\in\{-1,1\}$,
define
\begin{equation}
 c_s=\widehat y+sm\widehat x,\qquad
 v_s=\Sigma_{yy}+m^2\Sigma_{xx}+2sm\Sigma_{xy}.
 \label{exec22:hd-components}
\end{equation}
The projected set has endpoints
\begin{equation}
 L=\min_{s\in\{-1,1\}}(c_s-r\sqrt{v_s}),\qquad
 U=\max_{s\in\{-1,1\}}(c_s+r\sqrt{v_s}).
 \label{exec22:hd-endpoints}
\end{equation}

To verify the extrema, put $a_s=(sm,1)'$. Cauchy--Schwarz in the covariance
metric bounds $a_s'(\beta-\widehat\beta)$ by
$r\sqrt{a_s'\Sigma a_s}$, with equality at
$\beta=\widehat\beta\pm r\Sigma a_s/\sqrt{a_s'\Sigma a_s}$.
The identities $y+m|x|=\max_s(y+smx)$ and
$y-m|x|=\min_s(y+smx)$ then give
\eqref{exec22:hd-endpoints}. There is no missing interior interval:
the continuous map $(x,y,u)\mapsto y+um|x|$ sends the connected compact set
$E\times[-1,1]$ onto a connected compact subset of the real line.
On the event that the true $(x,y)$ lies in $E$, every effect allowed by its
relative-magnitude restriction lies in $[L,U]$. Thus this construction has
at least $1-\alpha$ coverage for each allowed target under
\eqref{exec22:hd-law}. This direct argument deliberately establishes a
conservative reference; it does not reproduce HonestDiD's optimized conditional,
hybrid or fixed-length procedures~\citep{exec22_honestdid_reference}.

Take $\widehat x=0$, $\widehat y=4$, $m=1$ and
\begin{equation}
 \Sigma=\begin{pmatrix}1&3/2\\3/2&4\end{pmatrix},
 \qquad \det\Sigma=7/4.
 \label{exec22:hd-example}
\end{equation}
The component variances are $2$ and $8$. At $\alpha=0.05$ the projection is
approximately $[-2.9233,10.9233]$. Dropping the off-diagonal covariance changes
both component variances to $5$ and produces $[-1.4733,9.4733]$ for a different
stipulated law. Plugging in a zero pretrend and using only the ordinary
postperiod interval gives approximately $[0.0801,7.9199]$; that positive lower
endpoint is not justified as an RM sensitivity bound. Under a separately
stipulated population example $x=0,y=4$, the target is $4$, and an offer cost
of USD $1$ per exposure gives net value $3$. Applying that same fixed cost
to the statistical projection gives approximately $[-3.9233,9.9233]$.
A claim of confidently positive incremental value therefore depends on which
uncertainty calculation is justified. No bank return is measured in this example
\citep{exec22_honestdid_reference}.

The paper's finite-normal and asymptotic statements also require their own
conditions. Its asymptotic size-control argument includes uniform normal
approximation, covariance eigenvalues bounded above and away from zero,
uniform covariance consistency and regularity of the restriction matrix.
The fixed-length interval's finite-sample near-optimality result has further
convexity and symmetry conditions; those do not apply automatically to relative
magnitude or added sign/shape restrictions. Neither a sensitivity curve nor
a convenient Gaussian approximation supplies these assumptions. Reporting a
favourable tuning value after inspecting the curve needs a separately justified
selection or simultaneous-inference rule~\citep{rambachan2023}.

\subsection{Defaults, grids and failed numerical searches}
\label{exec22:hd-numerics}

The pinned source exposes substantive distinctions between entrypoints.
The unrestricted smoothness interface defaults to a fixed-length confidence
interval and the first postperiod target. Its automatically selected tuning
grid uses a preperiod upper-bound calculation with significance level $0.05$
when multiple preperiods are present, even if the requested interval level
differs. The relative-magnitude interface defaults to the C-LF method,
ten tuning values between zero and two, and $1{,}000$ inversion grid points.
The \texttt{did} bridge has a default of $100$ grid points and forwards that
argument to relative magnitude but not to its smoothness call. A wrapper should
record the executed options explicitly and require deliberate tuning choices
for a decision analysis~\citep{exec22_honestdid_engine,exec22_honestdid_bridges}.

At the lower relative-magnitude level, membership is retained at grid nodes
after union across signed preperiod components. Default bounds are plus or
minus twenty post-target standard errors around zero. The inspected high-level
path reduces accepted nodes to their minimum and maximum. This is a numerical
summary, not a certificate that the continuous confidence set is connected,
fully searched or bounded. An empty set of accepted nodes on the coarse grid
$\{-1,1\}$ is compatible with the nonempty continuous set $[-0.2,0.2]$.
Accepted nodes $-2,-1,1,2$ with rejection at zero have hull $[-2,2]$, which
fills a rejected gap. Acceptance at a search boundary leaves the extent beyond
that boundary unresolved. These are elementary grid counterexamples, not
observed outputs from running the package~\citep{exec22_honestdid_engine,exec22_honestdid_reference}.

Solver status creates a different concern. In the inspected population-bound
helpers, when both linear-program statuses are nonzero, the smoothness path
returns the unadjusted target as both endpoints with a warning; the selected
relative-magnitude component path makes a corresponding return without a local
warning. The inspected alternative branch does not separately guard a single
failed optimum. These source paths do not prove that any particular runtime
input triggers a failure, but they rule out treating every returned point
interval as certified identification. A proposed adapter must distinguish
infeasibility, unboundedness, convergence failure and an unresolved status.
Capturing warnings alone is insufficient if a helper discards status; a pinned
patch or lower-level instrumentation may be required before admission
\citep{exec22_honestdid_engine}.

The same principle applies to covariance and restrictions. Nonconformable
dimensions cause refusal, while asymmetry and negative eigenvalues generate
warnings rather than a statistical repair. Fixed-length paths warn that
requested sign or shape restrictions are unused. The initial proposed contract
should refuse an unexplained invalid covariance and either refuse an ignored
restriction or report the executed, weaker restriction explicitly. A silent
positive-semidefinite projection would change the supplied uncertainty model;
it needs its own justification.

\subsection{Reuse choice and an implementable qualification path}
\label{exec22:hd-reuse}

The proposal retains the pinned R implementation as a reference and selects an
isolated R worker as a future integration candidate. Its complete repository
MIT licence permits reuse subject to notice retention and the stated disclaimer.
This does not qualify the licences, installation or numerical behavior of CVXR,
GLPK, ECOS and the other imported dependencies. The pinned maintenance notes
describe recent CVXR compatibility repairs, which make a frozen environment and
version-change review especially relevant. A new Python or C++ port would add
optimization and distributional parity obligations without a measured benefit.
The Stata implementation or another port would need its own source and interface
review~\citep{exec22_honestdid_engine}.

The conversion boundary is equally important. The selected \texttt{did} bridge
requires a dynamic aggregate and universal baseline, constructs covariance from
the influence matrix, requires consecutive event times, and removes the
normalised reference coefficient. This computation alone does not establish
that the influence representation has the right bank clustering or calendar
dependence. The selected \texttt{fixest} bridge uses internal model fields to
form an aggregation matrix $T$, then transforms coefficients and covariance as
$T\widehat\beta$ and $T\widehat\Sigma T'$. The worker must preserve cohort
ordering, target weights and the actual covariance selected upstream, and must
qualify compatibility with the chosen versions~\citep{exec22_honestdid_bridges}.

A bounded prototype should first reproduce population targets and option
forwarding on frozen contrasts; then exercise malformed order, singular inputs,
solver terminations, clipped grids and ignored restrictions. Coverage and
decision-value comparisons need prospectively specified generated panels,
sampling laws, tuning choices, failure denominators and resource budgets.
Ordinary DID, prespecified smoothness, relative magnitude and the conservative
projection supply useful matched comparators, with their differing assumptions
retained. Total cost includes event-study fitting, covariance construction,
conversion, repeated inversion, failure investigation and environment support.
Only that evidence can justify which R path to build and whether a later bank
experiment is worth its cost.

The rollout conclusion therefore has two obligations. Its reported interval
must retain the population, sampling law and method that produced it, and its
causal interpretation must withstand the defended counterfactual restrictions.
A reliable conversion cannot supply a missing trend assumption; a sensitivity
analysis cannot repair a lost covariance or a failed numerical bound. When
those obligations leave a wide or disconnected set, the unresolved result is
part of the decision. The next comparison examines how donor construction and
synthetic controls offer different counterfactual information, with their own
extrapolation and uncertainty costs.

